% Reproducibility Statement (to include near the end of the paper)
\section*{Reproducibility Statement}\label{sec:repro-appendix}
We scope claims to shallow MLPs on small, standard datasets and provide all code, presets, and artifacts to reproduce results on a laptop‑class machine.

\textbf{Datasets and splits.} MNIST, Fashion‑MNIST (subsets), 20 Newsgroups, AG News (small variant), UCR GunPoint, and California Housing. We use canonical train/test splits; validation subsets are formed deterministically with fixed seeds.

\textbf{Training protocol.} Seeds are fixed for initialization and data shuffling. We run $n\!\in\![2,5]$ repeats per configuration and report mean $\pm$ standard deviation. We sweep learning rate and ternary threshold $\tau$ within specified ranges. No early stopping is used.

\textbf{Implementation details.} Feedback strategies include backprop, DFA (float), ternary‑forward DFA, and structured feedback (orthonormal/Hadamard). Ternary quantization uses float shadow weights and deterministic or seeded stochastic rounding. Full configs are saved with each run.

\textbf{Environment.} CPU‑only runs on a host with Apple M3 (Darwin 24.6.0, arm64). Dependencies are pinned via \texttt{requirements‑lock.txt}. BLAS is provided via NumPy’s linked backend on macOS; all latencies are measured with high‑resolution timers on this host and reported alongside the device string in the Deployability table.

\textbf{Artifacts.} Per‑run metrics/manifests under \texttt{runs/}; aggregate tables and plots under \texttt{data/report/}. The Key Metrics table is auto‑generated for the paper build. Scripts: \texttt{reproduce\_all.sh}, \texttt{make report}, and \texttt{make paper}.

\textbf{Limitations.} DFA/FA often underperform on deep CNNs and large‑scale vision without adapted feedback; we do not claim results beyond our small MLPs.
